
\section{Results and Discussion}

\subsection{Expert Rule-Based Classification Logic}

Algorithm~2 presents the clinical decision rules formulated for rare disease detection, integrating hematological, growth, pulmonary, and neurological indicators. Each rule is built on codified ICD-10 diagnostic criteria, lab thresholds, and patient history, with a priority on early, interpretable identification of suspected Gaucher, ASMD, or Saposin C deficiency cases.

\begin{tcolorbox}[title=Algorithm 2: Expert Rules for Rare Disease Detection]
\textbf{Input:} Patient diagnosis, lab, clinical, and drug history \\
\textbf{Output:} Rule1 to Rule5 flags and final prediction

\begin{verbatim}
if (Organomegaly [R19.00, R19.05, R19.06, R19.09] OR Hepatosplenomegaly [R16.2] OR 
    Hepatomegaly [R16.0] OR Splenomegaly [R16.1]) AND Abnormal CBC (WBC, Hgb, 
    Platelet, Neutro Auto #) then
    Rule1 ← TRUE;
if Growth retardation [R62.50–R62.59] AND Rule1 then
    Rule2 ← TRUE;
if (Growth retardation OR Failure to thrive [R62.7] OR Height/Weight Centile < 3) AND 
    Rule1 AND Rule2 then
    Rule3 ← TRUE;
if Interstitial lung disease [J84.9, J84.10, J84.112, J84.114, J84.848] AND Rule1 then
    Rule4 ← TRUE;
if Rule1 AND (Hypotonia [P94.1, P94.2] OR Developmental delay [R62.50, R62.59]) then
    Rule5 ← TRUE;
positive_prediction ← Rule1 OR Rule2 OR Rule3 OR Rule4 OR Rule5;
\end{verbatim}
\end{tcolorbox}

\subsection{Performance Evaluation}

The expert rules were evaluated across three rare disease datasets using standard classification metrics (Table~\ref{tab:primary}) and extended metrics to address class imbalance (Table~\ref{tab:extended}).

\begin{table}[H]
\centering
\caption{Primary Classification Metrics}
\label{tab:primary}
\begin{tabular}{lccc}
\toprule
\textbf{Metric} & \textbf{ASMD (\%)} & \textbf{Gaucher (\%)} & \textbf{Saposin (\%)} \\
\midrule
Accuracy & 27.59 & 77.78 & 70.59 \\
Precision & 28.57 & 84.62 & 83.33 \\
Recall & 88.89 & 84.62 & 76.92 \\
Specificity & 0.00 & 60.00 & 50.00 \\
F1 Score & 43.24 & 84.62 & 80.00 \\
\bottomrule
\end{tabular}
\end{table}

\begin{table}[H]
\centering
\caption{Extended Evaluation Metrics}
\label{tab:extended}
\begin{tabular}{lccc}
\toprule
\textbf{Metric} & \textbf{ASMD (\%)} & \textbf{Gaucher (\%)} & \textbf{Saposin (\%)} \\
\midrule
Negative Predictive Value & 0.00 & 60.00 & 40.00 \\
Balanced Accuracy & 44.44 & 72.31 & 63.46 \\
\bottomrule
\end{tabular}
\end{table}

\subsection{Comparative Disease-Specific Outcomes}

\textbf{ASMD:} High sensitivity (Recall: 88.89\%) with specificity of 0\% and poor precision (28.57\%) due to nonspecific triggers.  
\textbf{Gaucher:} Consistently strong performance across all metrics, validating rule precision.  
\textbf{Saposin:} Moderate precision and recall with specificity at 50\%; misclassifications often stemmed from incomplete confirmatory records.

\subsection{Confusion Matrix Analysis}

\begin{table}[H]
\centering
\caption{Confusion Matrix for ASMD}
\begin{tabular}{lcc}
\toprule
\textbf{Actual \textbackslash{} Predicted} & Positive & Negative \\
\midrule
Positive & 8 (TP) & 1 (FN) \\
Negative & 20 (FP) & 0 (TN) \\
\bottomrule
\end{tabular}
\end{table}

\begin{table}[H]
\centering
\caption{Confusion Matrix for Gaucher}
\begin{tabular}{lcc}
\toprule
\textbf{Actual \textbackslash{} Predicted} & Positive & Negative \\
\midrule
Positive & 11 (TP) & 2 (FN) \\
Negative & 2 (FP) & 3 (TN) \\
\bottomrule
\end{tabular}
\end{table}

\begin{table}[H]
\centering
\caption{Confusion Matrix for Saposin}
\begin{tabular}{lcc}
\toprule
\textbf{Actual \textbackslash{} Predicted} & Positive & Negative \\
\midrule
Positive & 10 (TP) & 3 (FN) \\
Negative & 2 (FP) & 2 (TN) \\
\bottomrule
\end{tabular}
\end{table}

\subsection{Discussion and Interpretability}

The rule-based model is a transparent, clinically grounded diagnostic aid, ideal for early triage in low-resource settings.  

\textbf{Strengths:}  
\begin{itemize}
  \item High recall ensures minimal missed true cases.
  \item Rules mirror clinical logic, enhancing trust and adoption.
  \item Gaucher rules demonstrated strong discriminative power.
\end{itemize}

\textbf{Limitations:}  
\begin{itemize}
  \item ASMD rules prone to overflagging due to generic triggers.
  \item Static thresholds lack contextual adaptability.
  \item Limited standalone utility in ambiguous or overlapping cases.
\end{itemize}

\subsection{Future Integration}

These results support using expert rules as front-line screening tools in rare disease pipelines. Integrating them with machine learning classifiers or biochemical follow-ups may improve specificity while retaining clinical transparency.
